\section{Data analysis}
\label{sec:appendix-data-analysis}

This appendix reports the distributions behind the fingerprint. Every value is generated from the named analysis records, the windows evidence of the data inventory, the comparison record of the lunar target and the configuration, whose input digests are recorded; analyses whose records do not exist yet are named as such. All values here are Derived.

\paragraph{Site and region distributions.} Table~\ref{tab:distributions} summarises the
fingerprints of all \DistSiteReliefHCount{} windows of the lunar target site and all \DistHoluhraunReliefHCount{}
Holuhraun windows. The lunar windows have relief Hurst exponents from \DistSiteReliefHMin{} to
\DistSiteReliefHMax{} with a median of \DistSiteReliefHMedian, so the primary target window, at
\MoonTargetReliefH, is typical of its site. The Holuhraun windows span a much wider range, from
\DistHoluhraunReliefHMin{} to \DistHoluhraunReliefHMax, and their median relief $\nu$ at the
longest lag, \DistHoluhraunReliefNuLongMedian\,m, is more than an order of magnitude below the
target site median of \DistSiteReliefNuLongMedian\,m. The Hurst exponents of the target site windows lie
close to the highland median that \citet{rosenburg2011} report from LOLA, although our lags,
grid and estimator differ from theirs, so the agreement is a sanity check and not a
reproduction.

\begin{table}[h]
\centering
\caption{Derived fingerprint distributions per site, over the eligible windows of each site. $\nu$ values
are medians over windows, in metres, at the shortest and longest configured lag.}
\label{tab:distributions}
\footnotesize
\begin{tabular}{@{}llccccc@{}}
\toprule
Site & Variant & windows & $H$ min & $H$ median & $H$ max & median $\nu$ short, long \\
\midrule
Moon, target site & relief & \DistSiteReliefHCount{} & \DistSiteReliefHMin{} & \DistSiteReliefHMedian{} & \DistSiteReliefHMax{} & \DistSiteReliefNuShortMedian, \DistSiteReliefNuLongMedian{} \\
Moon, target site & texture & \DistSiteTextureHCount{} & \DistSiteTextureHMin{} & \DistSiteTextureHMedian{} & \DistSiteTextureHMax{} & \DistSiteTextureNuShortMedian, \DistSiteTextureNuLongMedian{} \\
Earth, Holuhraun & relief & \DistHoluhraunReliefHCount{} & \DistHoluhraunReliefHMin{} & \DistHoluhraunReliefHMedian{} & \DistHoluhraunReliefHMax{} & \DistHoluhraunReliefNuShortMedian, \DistHoluhraunReliefNuLongMedian{} \\
Earth, Holuhraun & texture & \DistHoluhraunTextureHCount{} & \DistHoluhraunTextureHMin{} & \DistHoluhraunTextureHMedian{} & \DistHoluhraunTextureHMax{} & \DistHoluhraunTextureNuShortMedian, \DistHoluhraunTextureNuLongMedian{} \\
\bottomrule
\end{tabular}
\end{table}

\paragraph{Bootstrap intervals and target percentiles.} Over the target site, the median relief
Hurst exponent is \SiteHMedian{} with an interquartile range from \SiteHQLow{} to \SiteHQHigh{}
and a moving block bootstrap interval for the median from \SiteHLow{} to \SiteHHigh{}
(\SiteBootReplicates{} replicates, blocks of two windows, a two dimensional extension of the block bootstrap of \citet{kunsch1989}, itself an extension of the bootstrap of \citet{efron1979}). The texture exponent has a median of \SiteTextureHMedian{} with an interval
from \SiteTextureHLow{} to \SiteTextureHHigh. Neighbouring windows overlap by half, so these
intervals describe one site, not the Moon. Within the site, the primary target window sits at the
\TargetHPercentile{}th percentile of relief $H$, the \TargetNuShortPercentile{}th percentile of
relief $\nu$ at the shortest lag and the \TargetNuLongPercentile{}th at the longest lag. Across the
error clones the percentile of $H$ stays between \TargetHCloneMin{} and \TargetHCloneMax.

\paragraph{Scale behaviour.} A deviogram of the whole site from \BreakMinLag\,m to
\BreakMaxLag\,m lags has a single slope Hurst exponent of \BreakSingleH. A hinged fit with a
break at \BreakM\,m is preferred over a straight line by a difference in the Bayesian information
criterion of \BreakDeltaBIC, with exponents of \BreakHBelow{} below and \BreakHAbove{} above the
break. \citet{rosenburg2011} report that a breakover, where present, typically sits near a
kilometre baseline on the Moon; our break lies in the same range, but the site is a single
massif, our lags come from a gridded and partly interpolated model rather than profiles, and the
comparison is a sanity check, not a reproduction. The fingerprint's lags all lie below this break.

\paragraph{Across bodies.} Table~\ref{tab:overlap} compares the three bodies in fingerprint space, using the windows that the search read: \BodyMoonWindows{} lunar windows of the target site, \BodyMarsWindows{} Martian windows from \BodyMarsGroups{} CTX models, and \BodyEarthWindows{} Earth windows from \BodyEarthGroups{} regions. The overlap coefficient is the summed minimum of the two normalised histograms on common bins, so it is one for identical distributions and zero for disjoint ones. The median relief Hurst exponent is \BodyMoonHMedian{} at the lunar site, \BodyMarsHMedian{} on Mars and \BodyEarthHMedian{} on Earth. The lunar site barely overlaps either other body at any lag, while the Martian and Earth samples overlap substantially. The samples are small, contiguous and not independent, so the table describes these windows and not the bodies.

\begin{table}[h]
\centering
\caption{Derived overlap coefficients between the bodies' fingerprint distributions over the windows read by the search. Short and long are relief $\ln \nu$ at the shortest and longest configured lag.}
\label{tab:overlap}
\footnotesize
\begin{tabular}{@{}lcccc@{}}
\toprule
Pair & relief $H$ & texture $H$ & relief $\ln \nu$ short & relief $\ln \nu$ long \\
\midrule
Earth and Moon & \OverlapEarthMoonReliefH{} & \OverlapEarthMoonTextureH{} & \OverlapEarthMoonNuShort{} & \OverlapEarthMoonNuLong{} \\
Mars and Moon & \OverlapMarsMoonReliefH{} & \OverlapMarsMoonTextureH{} & \OverlapMarsMoonNuShort{} & \OverlapMarsMoonNuLong{} \\
Earth and Mars & \OverlapEarthMarsReliefH{} & \OverlapEarthMarsTextureH{} & \OverlapEarthMarsNuShort{} & \OverlapEarthMarsNuLong{} \\
\bottomrule
\end{tabular}
\end{table}

\paragraph{Distances to the target.} Over the \ComparedCount{} compared Earth windows, relief
distances to the target run from \ReliefDistanceMin{} to \ReliefDistanceMax{} (median
\ReliefDistanceMedian) and texture distances from \TextureDistanceMin{} to \TextureDistanceMax{}
(median \TextureDistanceMedian). Because the two regions were chosen for radar availability and as
a terrain only control, these distances describe the evaluated windows and not the Earth as a
whole.

\paragraph{Analyses not yet available.} Roughness by land cover class and climate zone, the effect of a surface model mask, the declared regions against a global background sample, the comparison with published Martian roughness, and the overlap of the bodies in embedding space with and without gradient reversal have no records, so they are not reported.
